\documentclass[11pt,a4paper]{article}

\usepackage[utf8]{inputenc}
\usepackage[spanish,es-tabla]{babel}
\usepackage[T1]{fontenc}
\usepackage{lmodern}
\usepackage[margin=2.5cm]{geometry}
\usepackage{amsmath,amssymb}
\usepackage{booktabs}
\usepackage{tabularx}
\usepackage{longtable}
\usepackage{xcolor}
\usepackage{graphicx}
\usepackage{fancyhdr}
\usepackage{titlesec}
\usepackage{enumitem}
\usepackage[most]{tcolorbox}
\usepackage{listings}
\usepackage{microtype}
\usepackage[colorlinks=true,linkcolor=purple!80!black,urlcolor=blue!70!black,citecolor=purple!80!black]{hyperref}

% Paleta de colores corporativos bancarios
\definecolor{bankPrimary}{RGB}{76, 29, 149}    % Violeta institucional
\definecolor{bankSecondary}{RGB}{124, 58, 237} % Purpura brillante
\definecolor{bankDark}{RGB}{15, 23, 42}        % Slate oscuro
\definecolor{bankGray}{RGB}{241, 245, 249}     % Gris fondo
\definecolor{bankAccent}{RGB}{13, 148, 136}    % Teal / exito

% Encabezados y pies de pagina
\pagestyle{fancy}
\fancyhf{}
\fancyhead[L]{\small\color{gray}Sistema de Gestión de Créditos Bancarios · UMSA}
\fancyhead[R]{\small\color{gray}Manual de Usuario (INF-241)}
\fancyfoot[L]{\small\color{gray}Universidad Mayor de San Andrés · Carrera de Informática}
\fancyfoot[R]{\small\color{bankPrimary}\bfseries\thepage}
\renewcommand{\headrulewidth}{0.4pt}
\renewcommand{\footrulewidth}{0.4pt}

% Formato de titulos
\titleformat{\section}{\color{bankPrimary}\normalfont\Large\bfseries}{\thesection.}{1em}{}[\titlerule]
\titleformat{\subsection}{\color{bankSecondary}\normalfont\large\bfseries}{\thesubsection.}{1em}{}
\titleformat{\subsubsection}{\color{bankDark}\normalfont\normalsize\bfseries}{\thesubsubsection.}{1em}{}

% Cajas de alerta y notas
\newtcolorbox{bankBox}[2][]{
  colback=bankGray,
  colframe=bankPrimary,
  fonttitle=\bfseries\color{white},
  title=#2,
  rounded corners,
  boxrule=1pt,
  #1
}

\newtcolorbox{alertBox}[1]{
  colback=red!5!white,
  colframe=red!75!black,
  fonttitle=\bfseries\color{white},
  title={IMPORTANTE: #1},
  rounded corners,
  boxrule=1pt
}

\newtcolorbox{tipBox}[1]{
  colback=teal!5!white,
  colframe=bankAccent,
  fonttitle=\bfseries\color{white},
  title={RECOMENDACIÓN OPERATIVA: #1},
  rounded corners,
  boxrule=1pt
}

\begin{document}

% ==============================================================================
% PORTADA INSTITUCIONAL
% ==============================================================================
\begin{titlepage}
    \centering
    \vspace*{1cm}
    
    {\scshape\Large UNIVERSIDAD MAYOR DE SAN ANDRÉS\par}
    {\scshape\large FACULTAD DE CIENCIAS PURAS Y NATURALES\par}
    {\scshape\large CARRERA DE INFORMÁTICA\par}
    
    \vspace{1.5cm}
    \hrule height 2pt
    \vspace{0.4cm}
    {\Huge\bfseries\color{bankPrimary} MANUAL DE USUARIO Y GUÍA DE OPERACIÓN TÉCNICA\par}
    \vspace{0.4cm}
    {\Large\bfseries\color{bankSecondary} SISTEMA INTEGRAL DE GESTIÓN CREDITICIA\par}
    \vspace{0.4cm}
    \hrule height 1pt
    
    \vspace{1.5cm}
    {\large\bfseries MATERIA: INF-241 ANÁLISIS Y DISEÑO DE SISTEMAS I}\\[0.2cm]
    {\normalsize Paralelo ``B'' · Semestre II/2026}\par
    
    \vspace{2cm}
    \begin{minipage}{0.45\textwidth}
        \begin{flushleft} \large
            \textbf{Docente:}\\
            M. Sc. Patricia Sonia Trino Camacho
        \end{flushleft}
    \end{minipage}
    \hfill
    \begin{minipage}{0.45\textwidth}
        \begin{flushright} \large
            \textbf{Estudiante (Grupo 4):}\\
            Alex Arturo Jamachi Condori\\[0.1cm]
            \small\texttt{alexjamachicondori@gmail.com}
        \end{flushright}
    \end{minipage}

    \vfill
    {\large La Paz -- Bolivia\par}
    {\large Octubre de 2026\par}
\end{titlepage}

\newpage
\tableofcontents
\newpage

% ==============================================================================
% 1. INTRODUCCIÓN Y PROPÓSITO
% ==============================================================================
\section{Introducción y Alcance del Sistema}

El \textbf{Sistema de Gestión de Créditos para una Entidad Bancaria} es una plataforma de software transaccional, ergonómica y segura diseñada para administrar de punta a punta el ciclo operativo de intermediación financiera en colocación de activos de crédito. El sistema responde con rigor técnico a las directrices regulatorias de la \textbf{Autoridad de Supervisión del Sistema Financiero (ASFI)} de Bolivia y a la Ley N° 393 de Servicios Financieros.

\subsection{Objetivos del Manual}
\begin{itemize}[leftmargin=*]
    \item Servir como instrumento formal de capacitación para los operadores bancarios según su rol jerárquico.
    \item Detallar paso a paso los procedimientos operativos para cada caso de uso del sistema.
    \item Proporcionar directrices exactas sobre validaciones de riesgo, límites de atribución crediticia y cascada de cobro.
    \item Servir de manual de auditoría para la verificación de pistas inmutables y trazabilidad transaccional.
\end{itemize}

\subsection{Direcciones de Acceso al Entorno Operativo}
El sistema se encuentra desplegado y accesible a través de dos canales:
\begin{enumerate}
    \item \textbf{Servidor en la Nube (AWS Lightsail):} 
    \begin{itemize}
        \item Dirección IP Directa: \url{http://13.220.246.62}
        \item Canal Seguro SSL (Cloudflare): \url{https://recruitment-thing-ethics-reliable.trycloudflare.com}
    \end{itemize}
    \item \textbf{Entorno de Laboratorio / Local:} \url{http://localhost:5173}
\end{enumerate}

% ==============================================================================
% 2. SEGURIDAD Y ACCESO (RBAC)
% ==============================================================================
\section{Seguridad de Acceso y Gestión de Sesión (RNF01, RF16)}

El acceso al sistema está blindado mediante un modelo estricto de \textbf{Control de Acceso Basado en Roles (RBAC)} y autenticación con tokens criptográficos \textbf{JWT (JSON Web Tokens)}.

\subsection{Políticas de Seguridad en Pantalla}
\begin{itemize}[leftmargin=*]
    \item \textbf{Temporizador de Inactividad (RNF01):} Cada sesión cuenta con un tiempo máximo de vida de \textbf{15 minutos} (900 segundos). El temporizador regresivo se visualiza permanentemente en la barra superior del sistema. Al expirar el tiempo, el sistema cierra la sesión de forma inmediata y solicita volver a autenticarse para evitar accesos no autorizados.
    \item \textbf{Protección contra Fuerza Bruta:} Al acumularse \textbf{5 intentos fallidos consecutivos} de contraseña sobre una cuenta, el sistema bloquea preventivamente el usuario durante \textbf{30 minutos}, registrando la IP de origen en la bitácora de auditoría forense.
\end{itemize}

\subsection{Matriz de Usuarios y Credenciales Estandarizadas}
Para efectos de evaluación y pruebas piloto, todas las cuentas oficiales comparten la contraseña bancaria:
\begin{center}
    \colorbox{bankGray}{\texttt{\textbf{Bancario2026!}}}
\end{center}

\begin{table}[h!]
    \centering
    \small
    \begin{tabularx}{\textwidth}{llllX}
        \toprule
        \textbf{Rol RBAC} & \textbf{Login} & \textbf{Funcionario} & \textbf{Límite Atribución} & \textbf{Módulos Autorizados} \\
        \midrule
        \textbf{OFICIAL} & \texttt{atorrez} & Alvaro Torrez & Bs 0,00 & Clientes (P10), Solicitudes (P05). \\
        \textbf{ANALISTA} & \texttt{lrojas} & Laura Rojas & Bs 0,00 & Solicitudes, Consulta Buró (P06), Garantías (P08). \\
        \textbf{JEFE} & \texttt{cmendoza} & Carlos Mendoza & Bs 70.000,00 & Bandeja Aprobación (P02), Desembolsos (CU11), Reportes. \\
        \textbf{CAJERO} & \texttt{cajero1} & Mario Fernandez & Bs 0,00 & Ventanilla de Caja y Pagos (CU12/CU13), Recibos. \\
        \textbf{GERENTE} & \texttt{gerente1} & Sofia Valdivia & Sin Límite & Aprobación montos mayores (> Bs 70.000), Reportes Cartera. \\
        \textbf{ADMINISTRADOR} & \texttt{admin} & Administrador & Total & Usuarios, Roles, Desbloqueo, Auditoría Forense. \\
        \bottomrule
    \end{tabularx}
    \caption{Matriz oficial de asignación de roles, atribuciones y credenciales del sistema.}
\end{table}

% ==============================================================================
% 3. MÓDULO OFICIAL DE CRÉDITO
% ==============================================================================
\section{Módulo del Oficial de Crédito (\texttt{atorrez})}

El Oficial de Crédito es el responsable de la prospección, recepción documental, registro del cliente y conformación de la carpeta crediticia inicial.

\subsection{Registro y Mantenimiento de Clientes (RF01, RF02, P10)}
\begin{enumerate}
    \item Inicie sesión con el usuario \texttt{atorrez}.
    \item En la barra de navegación superior, haga clic en el botón \textbf{Clientes (P10)}.
    \item Para ingresar un cliente nuevo, presione el botón \textbf{+ Nuevo Cliente}.
    \item Complete los campos obligatorios del formulario modal:
    \begin{itemize}
        \item \textbf{Cédula de Identidad (CI):} Debe ser única en el sistema bancario. Si el carnet ya existe, el sistema emitirá un bloqueo impidiendo duplicidades.
        \item \textbf{Nombres y Apellidos Completos.}
        \item \textbf{Dirección de Domicilio y Teléfono de Contacto.}
        \item \textbf{Ingresos Mensuales Declarados (Bs):} Importe neto respaldado documentalmente.
        \item \textbf{Egresos Mensuales Declarados (Bs):} Gastos fijos de subsistencia y deudas externas.
    \end{itemize}
    \item Haga clic en \textbf{Guardar Cliente}.
\end{enumerate}

\subsection{Registro de Solicitudes de Crédito (RF03, P05)}
\begin{enumerate}
    \item Seleccione la pestaña \textbf{Solicitudes de Crédito} y presione \textbf{+ Nueva Solicitud}.
    \item Seleccione el cliente titular previamente creado en la lista desplegable.
    \item Elija el \textbf{Tipo de Crédito}:
    \begin{itemize}
        \item \textbf{CONSUMO:} Tasa del $15.00\%$ anual (Plazos de 12 a 60 meses).
        \item \textbf{VEHICULAR:} Tasa del $10.00\%$ anual (Plazos de 12 a 72 meses).
        \item \textbf{VIVIENDA:} Tasa del $7.50\%$ anual (Plazos de 60 a 240 meses).
        \item \textbf{MICROCRÉDITO:} Tasa del $18.00\%$ anual (Plazos de 6 a 36 meses).
    \end{itemize}
    \item Ingrese el \textbf{Monto Solicitado (Bs)} y el \textbf{Plazo en Meses}.
    \item \textbf{Simulación Instantánea:} Observe que el sistema calcula automáticamente en tiempo real la \textbf{Cuota Mensual Estimada} bajo Sistema Francés y el indicador preliminar de relación cuota/ingreso (DTI).
    \item Ingrese el destino de los fondos y presione \textbf{Registrar Solicitud}. La solicitud quedará en estado inicial \texttt{REGISTRADA}.
\end{enumerate}

% ==============================================================================
% 4. MÓDULO ANALISTA DE RIESGOS
% ==============================================================================
\section{Módulo del Analista de Riesgos (\texttt{lrojas})}

El Analista de Riesgos verifica la solvencia moral y financiera del prestatario, consulta las centrales de riesgo y valida el respaldo de garantías reales o personales.

\subsection{Consulta al Buró de Información Crediticia ASFI (RF04, P06)}
\begin{enumerate}
    \item Inicie sesión con el usuario \texttt{lrojas} y diríjase al listado de solicitudes.
    \item En la solicitud que se encuentre en estado \texttt{REGISTRADA}, presione el botón \textbf{Evaluar Solicitud}.
    \item En el modal de evaluación, presione el botón \textbf{Consultar Buró de Riesgo}.
    \item El sistema se comunicará con el adaptador simulado del Buró ASFI, devolviendo:
    \begin{itemize}
        \item \textbf{Score Crediticio (0 a 1000 puntos):} Valores mayores o iguales a 600 califican como solventes.
        \item \textbf{Calificación Histórica:} \texttt{SIN MORA (A)}, \texttt{MORA LEVE (B)} o \texttt{MORA CRÍTICA (C/D)}.
    \end{itemize}
\end{enumerate}

\subsection{Evaluación Financiera y Ratio DTI (RF05, RN02)}
El analista introduce los ingresos y egresos formalmente verificados tras la inspección documental:
\begin{equation}
    DTI = \frac{\text{Cuota Mensual Estimada}}{\text{Ingresos Verificados} - \text{Egresos Verificados}} \times 100
\end{equation}
\begin{alertBox}{Regla Normativa ASFI (RN02)}
    El ratio de endeudamiento (\textbf{DTI}) no puede superar el \textbf{40.00\%}. Si el resultado supera este umbral, el sistema marcará automáticamente la solicitud como con riesgo excesivo.
\end{alertBox}

\subsection{Constitución y Cobertura de Garantías (RF07, RN03, P08)}
\begin{enumerate}
    \item En la tabla de la solicitud, presione el botón \textbf{Garantías (P08)}.
    \item Registre el tipo de bien (\texttt{PERSONAL}, \texttt{HIPOTECARIA} o \texttt{PRENDARIA}) y el valor comercial tasado.
    \item \textbf{Regla del 120\% (RN03):} El sistema calcula automáticamente la cobertura porcentual:
    \begin{equation}
        \text{Cobertura} = \frac{\sum \text{Valor Comercial Garantías}}{\text{Monto Solicitado}} \times 100 \ge 120\%
    \end{equation}
    \item Si la cobertura es inferior al $120\%$, no se podrá emitir dictamen favorable.
\end{enumerate}

Una vez cumplidos los requisitos, el analista emite el dictamen formal: \texttt{FAVORABLE}, \texttt{OBSERVADO} o \texttt{RECHAZADO}. La solicitud pasa a estado \texttt{EN\_EVALUACION}.

% ==============================================================================
% 5. MÓDULO JEFE DE CRÉDITOS Y DESEMBOLSOS
% ==============================================================================
\section{Módulo del Jefe de Créditos y Desembolsos (\texttt{cmendoza})}

El Jefe de Créditos ejerce la potestad de resolución y formalización jurídica del activo de crédito.

\subsection{Bandeja de Aprobación y Límites de Atribución (RF08, RN06, P02)}
\begin{enumerate}
    \item Inicie sesión con el usuario \texttt{cmendoza}.
    \item Acceda a la \textbf{Bandeja Aprobación (P02)}. Se listarán exclusivamente las solicitudes con dictamen emitido.
    \item Presione el botón \textbf{Aprobar}.
    \item \textbf{Control de Atribución Operativa (RN06):} El sistema verifica que el monto solicitado no sobrepase los \textbf{Bs 70.000,00} (atribución asignada al Jefe). Para montos superiores, el sistema restringe la operación e instruye su escalamiento a Gerencia General.
    \item Confirme la aprobación. El sistema transiciona la solicitud a \texttt{APROBADA} y crea el registro de crédito en estado \texttt{PENDIENTE\_DESEMBOLSO}.
\end{enumerate}

\subsection{Contrato Legal y Plan de Amortización Francés (RF10, RF11)}
Tras la aprobación se habilitan inmediatamente dos botones de auditoría:
\begin{itemize}
    \item \textbf{Ver Contrato Legal (RF10):} Despliega el documento oficial correlativo (\texttt{CT-2026-XXXXX}) con los datos del acreedor bancario, titular, C.I., capital prestado, tasa compensatoria, tasa penal y condiciones resolutivas.
    \item \textbf{Ver Cronograma Francés (RF11):} Presenta la tabla completa de cuotas, desglosando período a período: Fecha de Vencimiento, Saldo Deudor Inicial, Cuota Fija, Capital Amortizado, Interés Compensatorio y Saldo Restante. Compruebe que la última cuota ajusta con exactitud los centavos de redondeo dejando el saldo en \textbf{Bs 0,00}.
\end{itemize}

\subsection{Desembolso en Core Bancario (RF12, CU11)}
\begin{enumerate}
    \item En la barra de navegación, ingrese a \textbf{Desembolsos Core}.
    \item Localice el crédito formalizado en estado \texttt{PENDIENTE\_DESEMBOLSO}.
    \item Presione el botón \textbf{Desembolsar Fondos}.
    \item Confirme el diálogo del navegador. El sistema se comunica atómicamente con el Core Bancario, emitiendo el código de transacción único (ej. \texttt{CORE-TXN-307003}) y cambiando el estado del crédito a \texttt{VIGENTE}.
\end{enumerate}

% ==============================================================================
% 6. MÓDULO CAJERO DE VENTANILLA
% ==============================================================================
\section{Módulo de Ventanilla y Cobranza en Caja (\texttt{cajero1})}

El Cajero es el encargado de procesar los reembolsos de cartera, imputando los abonos bajo la normativa contable de prelación de pagos.

\subsection{Búsqueda y Cobro de Cuotas (RF13, CU12)}
\begin{enumerate}
    \item Inicie sesión con el usuario \texttt{cajero1} e ingrese a \textbf{Caja y Cobranza}.
    \item En el panel izquierdo, seleccione el crédito del cliente por su nombre o documento de identidad.
    \item En el panel derecho se desplegará el saldo de capital vivo y la lista cronológica de cuotas pendientes.
    \item En la primera cuota vencida o próxima a vencer, presione \textbf{Cobrar Cuota}.
\end{enumerate}

\subsection{Cascada Contable de Imputación de Pagos (RN10)}
Cuando se registra el abono, el motor bancario aplica los fondos en el siguiente orden estricto:
\begin{center}
    \textbf{1° Interés Penal por Mora} $\longrightarrow$ \textbf{2° Interés Compensatorio} $\longrightarrow$ \textbf{3° Amortización a Capital}
\end{center}
\begin{tipBox}{Emisión de Comprobante Fiscal (RF14)}
    Al procesar el cobro, el sistema genera de forma inmutable el \textbf{Comprobante Oficial de Caja} con código correlativo (\texttt{REC-XXXXX}), detallando el desglose exacto de los importes imputados, el cajero que procesó la transacción y el nuevo saldo deudor del cliente.
\end{tipBox}

\subsection{Control Diario de Mora (CU13)}
El cajero cuenta con el botón \textbf{Ejecutar Control Diario de Mora}. Esta rutina evalúa las cuotas vencidas y actualiza el interés penal devengado día a día sobre el capital en mora a la tasa anual reglamentada.

% ==============================================================================
% 7. MÓDULOS GERENCIALES Y ADMINISTRATIVOS
% ==============================================================================
\section{Módulos de Gerencia y Administración}

\subsection{Gerencia de Sucursal (\texttt{gerente1}): Reportes de Cartera (RF15)}
El Gerente cuenta con el panel de \textbf{Reportes Ejecutivos de Cartera}, donde visualiza en tiempo real:
\begin{itemize}
    \item Total de créditos formalizados y colocados en Bolivianos.
    \item Saldo de cartera vigente saludable.
    \item Saldo de cartera en mora y porcentaje de morosidad institucional ($0.00\%$).
    \item Total acumulado de cobranzas procesadas en caja.
\end{itemize}

\subsection{Administrador del Sistema (\texttt{admin}): Usuarios y Auditoría (RF16, RNF06)}
\begin{itemize}
    \item \textbf{Gestión de Usuarios:} Creación de nuevos operadores bancarios, asignación de roles RBAC, configuración de límites de atribución operativa y activación/desactivación de cuentas.
    \item \textbf{Pistas de Auditoría Forense (RN12):} Módulo inalterable que audita en tiempo real cada transacción del sistema:
    \begin{itemize}
        \item Marca temporal con precisión de milisegundos.
        \item Usuario responsable y su rol.
        \item Operación ejecutada (\texttt{LOGIN}, \texttt{CREAR\_CLIENTE}, \texttt{REGISTRAR\_SOLICITUD}, \texttt{APROBAR\_CREDITO}, \texttt{DESEMBOLSAR}, \texttt{REGISTRAR\_PAGO}).
        \item Dirección IP del cliente y resultado de la operación (\texttt{EXITOSO} / \texttt{FALLIDO}).
    \end{itemize}
\end{itemize}

% ==============================================================================
% 8. PREGUNTAS FRECUENTES Y RESOLUCIÓN DE INCIDENCIAS
% ==============================================================================
\section{Preguntas Frecuentes y Solución de Incidencias}

\begin{description}
    \item[¿Qué ocurre si la sesión se cierra automáticamente?]  
    Por normativa de seguridad bancaria (RNF01), el sistema cierra la sesión tras 15 minutos sin interacción. Simplemente vuelva a introducir su login y contraseña.
    
    \item[¿Por qué un Jefe no puede aprobar un crédito de Bs 80.000?]  
    Porque el límite de atribución operativa del rol Jefe de Créditos es de Bs 70.000,00 (RN06). Todo crédito superior debe ser aprobado por el usuario con rol \texttt{GERENTE}.
    
    \item[¿Por qué no se puede desembolsar un crédito?]  
    Verifique que el crédito se encuentre en estado \texttt{PENDIENTE\_DESEMBOLSO}. Un crédito rechazado o aún en evaluación no puede recibir fondos del Core Bancario (RN09).
\end{description}

\vspace{1cm}
\begin{center}
    \hrulefill\\
    \small\textbf{Sistema de Gestión de Créditos Bancarios · INF-241 UMSA}\\
    \small Cifrado BCrypt · Sesión JWT 15 min · Trazabilidad Forense Inmutable · Normativa ASFI
\end{center}

\end{document}
